\section{Hilbert systems, tableaux, and cut-free sequents}\label{sec:systems}

\subsection{A complete Hilbert presentation with explicit proof translation}

Consider the following nine axiom schemes, together with modus ponens.
Every displayed letter can be replaced by an arbitrary formula:
\[
\begin{array}{ll}
 \mathrm K&A\to(B\to A),\\
 \mathrm S&(A\to(B\to C))\to((A\to B)\to(A\to C)),\\
 \wedge I&A\to(B\to A\wedge B),\\
 \wedge E_1&(A\wedge B)\to A,\\
 \wedge E_2&(A\wedge B)\to B,\\
 \vee I_1&A\to A\vee B,\\
 \vee I_2&B\to A\vee B,\\
 \vee E&(A\to C)\to((B\to C)\to((A\vee B)\to C)),\\
 \bot E&\bot\to A.
\end{array}
\]
Call the calculus $H_I$. Add the scheme $A\vee\neg A$ to obtain $H_C$.
A Hilbert proof is a finite list, each entry annotated as an assumption, a
scheme instance with its substitutions, or modus ponens from earlier indices.
The checker is primitive recursive. Axiom schemes make a separate rule of
uniform substitution unnecessary, though Theorem~\ref{thm:structural} gives
that rule as a proof transformation.

\begin{lemma}[Effective Hilbert deduction theorem]\label{thm:hilbert-deduction}
From a length-$l$ Hilbert proof of $B$ from $\Gamma,A$, a primitive recursive
transformation produces a proof of $A\to B$ from $\Gamma$, using at most
$5l$ lines before optional sharing. It applies to both $H_I$ and $H_C$.
\end{lemma}
\begin{proof}
Process each original line $D$, producing a block ending in $A\to D$.
If $D$ is an axiom or an assumption from $\Gamma$, copy $D$, use
$D\to(A\to D)$, and apply modus ponens. If $D=A$, use the following five
lines:
\[
\begin{gathered}
 A\to((A\to A)\to A),\qquad A\to(A\to A),\\
 (A\to((A\to A)\to A))\to((A\to(A\to A))\to(A\to A)),\\
 (A\to(A\to A))\to(A\to A),\qquad A\to A.
\end{gathered}
\]
The first two are instances of K, the third of S, and the last two are
modus ponens. If $D$ was obtained from $E$ and $E\to D$, earlier blocks
provide $A\to E$ and $A\to(E\to D)$; an S instance and two modus ponens
steps give $A\to D$. This is at most five lines per original line.
Index translation and checking the resulting list are bounded scans.
\end{proof}

\begin{theorem}[Hilbert--natural-deduction equivalence]\label{thm:hilbert-equivalence}
There are primitive recursive proof translations in both directions between
$H_I$ and \IPC, and between $H_C$ and \CPC, preserving finite contexts
and conclusions. Both Hilbert systems are sound, and $H_C$ is Boolean-complete.
\end{theorem}
\begin{proof}
Each Hilbert scheme has a fixed natural-deduction proof template. K assumes
$A$ and then $B$ and returns $A$. S assumes functions $A\to(B\to C)$ and
$A\to B$, then assumes $A$, applies both functions, and obtains $C$.
The conjunction and disjunction schemes are the corresponding introduction
or elimination rules with their premises discharged. The falsehood scheme
uses falsehood elimination. Modus ponens is implication elimination, and
the extra classical scheme is an axiom of \CPC.

Conversely, translate a natural-deduction tree recursively. Assumptions
remain assumptions. Introduction and elimination of conjunction and
disjunction introduction use the corresponding scheme with modus ponens.
Implication elimination already is modus ponens. Implication introduction
uses the effective deduction theorem. For disjunction elimination, translate
its three premises; apply the deduction theorem to the two branch proofs,
obtaining $A\to C$ and $B\to C$; combine them and the proof of $A\vee B$
using the $\vee E$ scheme. Falsehood elimination uses $\bot E$.
Each transformation has an explicit primitive recursion on the original
tree; repeated use of the linear deduction transformation gives, for
example, a bound $100^{l+1}$ on the output line count for a tree of $l$
rule nodes. Formula-label lengths add only specified coding operations.
Soundness and classical completeness follow by composition with
Theorems~\ref{thm:bool-sound} and~\ref{thm:cpc-complete}.
\end{proof}

\subsection{The classical three-scheme presentation}

The root article displays the schemes K, S, and
\[
 \mathrm{L3}:\quad(\neg A\to\neg B)\to(B\to A).
\]
We can prove completeness for their following precisely specified version.
Take the language $\{p_i,\bot,\to\}$ and the abbreviation
$\neg A=A\to\bot$. This is a definitional version of the displayed
\L{}ukasiewicz schemes. The historical system with a separate primitive
negation symbol has a different literal grammar; our theorem refers to
the declared grammar and does not silently identify those syntaxes.

\begin{theorem}[Three-scheme completeness in the declared language]\label{thm:three-axiom}
K, S, L3, and modus ponens are sound and complete for Boolean truth in
$\{\bot,\to\}$. Proof translation and proof construction are primitive
recursive. Classical conjunction and disjunction can be added by the
definitions
\[
 A\wedge_c B:=\neg(A\to\neg B),\qquad
 A\vee_c B:=\neg A\to B.
\]
\end{theorem}
\begin{proof}
The deduction theorem used only K and S and therefore still applies.
It gives all purely implicational natural deductions. In particular,
$X\to\neg\neg X$ is derivable: assume $X$, assume $X\to\bot$, apply the
second assumption to the first, and discharge both in reverse order.
Also $\bot\to\bot$ is derivable. By K,
$\neg A\to(\bot\to\bot)$ is derivable. L3 with $B=\bot$ now yields
$\bot\to A$, so falsehood elimination has been recovered.

Use $X=\neg A$ in double-negation introduction to obtain
$\neg A\to\neg\neg\neg A$. L3 with $B=\neg\neg A$ yields
$\neg\neg A\to A$. Thus double-negation elimination is recovered as well.
To recover classical case analysis, suppose the two branch proofs give
$A\to C$ and $\neg A\to C$. Under an extra assumption $\neg C$, the first
branch implies $\neg A$; the second then gives $C$, contradicting $\neg C$.
Discharge it and use double-negation elimination to conclude $C$.
The signed evaluation construction for formulas in this smaller language
uses only implication rules and falsehood elimination; its merging pass
uses exactly this derived case rule. It therefore supplies completeness.
Soundness is the two-bit truth-table verification of K, S, and L3, followed
by modus ponens. Finally the two definitions have the asserted Boolean
tables, so recursive expansion gives the announced classical extensions.
All derivations just described have fixed templates, and the earlier
proof constructors supply the uniform bounds.
\end{proof}

The six schemes attributed to Frege in the source article are also
certifiable in this definitional language. Besides K and S, the exchange
scheme $(D\to(B\to A))\to(B\to(D\to A))$ is proved by assuming its three
antecedents and applying them in the indicated order. Contraposition
$(B\to A)\to(\neg A\to\neg B)$ assumes $B$ to obtain a contradiction.
Double-negation introduction was just proved. The remaining scheme is
double-negation elimination. Conversely K, S, and that last scheme already
give falsehood elimination: from $\bot$ obtain $\neg\neg A$ by weakening
under $\neg A$, then eliminate the double negation. They then give the
same case-analysis completeness proof. This proves the corresponding
six-scheme version sound and complete, without relying on historical
attribution to establish any mathematical step.

\subsection{Analytic tableaux with explicit fuel}

A signed formula is $\mathsf T A$ or $\mathsf F A$. It is satisfied by a
Boolean row precisely when $A$ has the indicated value. A branch is a
finite collection of signed formulas. Expanding a compound formula removes
that occurrence and adds the entries in the following table. A vertical
bar denotes two alternative branches:
\[
\begin{array}{c|c}
\text{principal formula}&\text{replacement}\\\hline
\mathsf T(A\wedge B)&\mathsf T A,\mathsf T B\\
\mathsf F(A\wedge B)&\mathsf F A\mid\mathsf F B\\
\mathsf T(A\vee B)&\mathsf T A\mid\mathsf T B\\
\mathsf F(A\vee B)&\mathsf F A,\mathsf F B\\
\mathsf T(A\to B)&\mathsf F A\mid\mathsf T B\\
\mathsf F(A\to B)&\mathsf T A,\mathsf F B.
\end{array}
\]
A branch closes on $\mathsf T\bot$ or on both $\mathsf T p_i,\mathsf F p_i$.
An occurrence $\mathsf F\bot$ may be removed. Fix a deterministic order for
choosing the next unexpanded compound formula. The total number of its
primitive nodes, summed over pending formulas, strictly decreases along
every branch: a nonbranching expansion removes one root node, and a
branching expansion retains only a proper subformula on each child.

\begin{theorem}[Tableau soundness, completeness, and termination]\label{thm:tableaux}
For an input branch with total node count $L$, the expansion process has
depth at most $L$ and fewer than $2^{L+1}$ nodes. It returns either a closed
tableau or a satisfying valuation. A branch has a closed tableau if and
only if it is unsatisfiable. All assertions and the proof-or-model
construction are provable in \ipra.
\end{theorem}
\begin{proof}
Each table entry is a Boolean equivalence: satisfying the parent is
equivalent to satisfying its only child, or at least one of its two
children. For example an implication is false exactly when its antecedent
is true and its consequent false, and is true exactly when the antecedent
is false or the consequent true. The decreasing node measure gives the
displayed finite bound by induction. Use that bound as recursion fuel,
padding completed branches with idle steps if a fixed loop is desired.

At a saturated open leaf all remaining constraints are consistent atomic
signs. Assign one to atoms with a positive sign and zero to the others.
No negatively constrained atom receives one because the branch is open.
Every sign at the leaf is therefore satisfied. The local equivalences
transport this valuation back up the chosen path to the original branch.
If every leaf closes, no valuation satisfies a leaf, and the same local
equivalences show that no valuation satisfies the root. The algorithm
decides which of these finite outcomes occurs and chooses the first open
leaf in the second case. This proves both directions constructively.
\end{proof}

To decide $\Gamma\models_b A$, start with $\mathsf T B$ for each
$B\in\Gamma$ and $\mathsf F A$. The theorem makes tableaux equivalent to
the truth-table consequence test, and hence to \CPC\ derivability.

\subsection{A classical cut-elimination algorithm by semantic reconstruction}

A multiple-conclusion sequent $\Gamma\Rightarrow\Delta$ says that every
valuation satisfying all members of $\Gamma$ satisfies some member of
$\Delta$. The empty succedent says there is no such valuation. Use initial
sequents $\Gamma,A\Rightarrow\Delta,A$ and
$\Gamma,\bot\Rightarrow\Delta$, weakening, exchange and contraction,
and the following rules:
\[
\begin{gathered}
 \dfrac{\Gamma,A,B\Rightarrow\Delta}{\Gamma,A\wedge B\Rightarrow\Delta},\qquad
 \dfrac{\Gamma\Rightarrow\Delta,A\quad\Gamma\Rightarrow\Delta,B}
       {\Gamma\Rightarrow\Delta,A\wedge B},\\[5pt]
 \dfrac{\Gamma,A\Rightarrow\Delta\quad\Gamma,B\Rightarrow\Delta}
       {\Gamma,A\vee B\Rightarrow\Delta},\qquad
 \dfrac{\Gamma\Rightarrow\Delta,A,B}{\Gamma\Rightarrow\Delta,A\vee B},\\[5pt]
 \dfrac{\Gamma\Rightarrow\Delta,A\quad\Gamma,B\Rightarrow\Delta}
       {\Gamma,A\to B\Rightarrow\Delta},\qquad
 \dfrac{\Gamma,A\Rightarrow\Delta,B}{\Gamma\Rightarrow\Delta,A\to B}.
\end{gathered}
\]
The cut rule, if included, is
\[
 \dfrac{\Gamma\Rightarrow\Delta,A\quad\Gamma,A\Rightarrow\Delta}
       {\Gamma\Rightarrow\Delta}.
\]
The shared-context formulation suffices; weakening gives the usual rule
with different contexts.

\begin{theorem}[Analytic completeness and effective cut elimination]\label{thm:cut}
Every Boolean-valid finite sequent has a cut-free proof in this calculus
using only subformulas of its end sequent. A primitive recursive function
maps any checked proof, including one with cuts, to a cut-free proof of the
same sequent. For end-sequent node count $L$, a bound
$10(L+1)2^{L+1}$ on rule nodes is sufficient with explicit structural steps.
\end{theorem}
\begin{proof}
Read a sequent as the tableau branch with left formulas true and right
formulas false. Each tableau expansion, read from child sequents to the
parent, is exactly one of the six displayed rules. The closed leaves are
initial sequents. Removing $\mathsf F\bot$ in search becomes right weakening
by $\bot$ in the resulting proof. An atom may be marked as processed for
the search without changing its sequent; it remains available for leaf
closure. Thus a closed tableau translates directly into a cut-free proof.
Move a chosen principal formula to the end of its side with at most $L$
adjacent exchanges, insert its children there, and leave the other entries
in their previous order. The tableau bound, together with this linear
number of ordering and weakening steps at a node, gives the stated generous bound. All formulas
encountered are subformulas of the end sequent.

For an input proof with cuts, first use its finite Boolean soundness
invariant. The displayed rules are sound by their truth tables. Cut is
sound because on a row satisfying $\Gamma$ and falsifying $\Delta$, its
first premise forces $A$ true and its second forbids $A$ true. Soundness
therefore proves validity of the end sequent, so its tableau cannot have
an open leaf and yields the required cut-free proof. This is a
\emph{reconstruction} algorithm, not a claim about the termination rate of
every possible local Gentzen reduction strategy. It proves cut elimination
as the existence of a primitive recursive transformation with a concrete
correctness contract.
\end{proof}

\section{Elementary intuitionistic laws and proof terms}\label{sec:laws}

\subsection{A library of constructive laws}

The following identities mean that both implications have \IPC\ proofs.
The list is schematic and thus yields infinitely many instances:
\[
\begin{array}{ll}
 A\wedge B\leftrightarrow B\wedge A,&A\vee B\leftrightarrow B\vee A,\\
 (A\wedge B)\wedge C\leftrightarrow A\wedge(B\wedge C),&
 (A\vee B)\vee C\leftrightarrow A\vee(B\vee C),\\
 A\wedge A\leftrightarrow A,&A\vee A\leftrightarrow A,\\
 A\wedge\top\leftrightarrow A,&A\vee\bot\leftrightarrow A,\\
 A\wedge\bot\leftrightarrow\bot,&A\vee\top\leftrightarrow\top,\\
 A\wedge(A\vee B)\leftrightarrow A,&A\vee(A\wedge B)\leftrightarrow A,\\
 A\wedge(B\vee C)\leftrightarrow(A\wedge B)\vee(A\wedge C),&\phantom{X}\\
 A\vee(B\wedge C)\leftrightarrow(A\vee B)\wedge(A\vee C).&
\end{array}
\]
Further constructive equivalences include
\[
\begin{split}
 (A\wedge B)\to C&\ \leftrightarrow\ A\to(B\to C),\\
 A\to(B\wedge C)&\ \leftrightarrow\ (A\to B)\wedge(A\to C),\\
 (A\vee B)\to C&\ \leftrightarrow\ (A\to C)\wedge(B\to C),\\
 \neg(A\vee B)&\ \leftrightarrow\ \neg A\wedge\neg B,\\
 \neg(A\wedge B)&\ \leftrightarrow\ A\to\neg B,\\
 \neg\neg\neg A&\ \leftrightarrow\ \neg A.
\end{split}
\]

\begin{theorem}[Constructive law certificates]\label{thm:int-laws}
Every equivalence in the two displays has a fixed \IPC\ proof template.
So do modus ponens, hypothetical syllogism, contraposition, disjunctive
syllogism, and constructive dilemma. Substitution gives primitive
recursive proof constructors for every instance.
\end{theorem}
\begin{proof}
For conjunction commutativity and associativity, project the required
components and introduce them in the target order; reversing that order
gives the converse. Idempotence projects one copy in one direction and
duplicates an available proof in the other. A proof of $\top$ is the
identity proof of $\bot\to\bot$. Thus conjunction with $\top$ may be
introduced or projected away. Conjunction with $\bot$ yields $\bot$ by
projection, and $\bot$ yields both components by elimination.

For disjunction commutativity use two branches, introducing the other
injection in each. For associativity use three branches, introducing the
same alternative with its new parenthesization; the reverse uses the
opposite parenthesization. For idempotence both branches return $A$.
In $A\vee\bot$ the $A$ branch returns $A$ and the falsehood branch returns
$A$ by elimination; its converse is the left injection. Both directions
of $A\vee\top\leftrightarrow\top$ follow by the always available proof of
$\top$ and by its right injection. For absorption, projection gives
$A$ from $A\wedge(A\vee B)$, and a given $A$ supplies that conjunction.
For $A\vee(A\wedge B)$, either branch supplies $A$; conversely inject $A$.

For the first distributivity law, project $A$ and split $B\vee C$; pair
$A$ with the chosen disjunct and inject the resulting conjunction. Conversely
split the outer disjunction, project its components, introduce $B\vee C$,
and pair with $A$. For the second law, a proof of $A$ gives both disjunctions
and a proof of $B\wedge C$ gives their right injections. Conversely split
$A\vee B$: its $A$ branch already gives the result; in its $B$ branch split
$A\vee C$, returning $A$ in one case and $B\wedge C$ in the other.

Currying uses an assumption of $A$ followed by one of $B$ to form their
pair; uncurrying projects a given pair and applies the curried function.
For implication into a conjunction, apply the function and project, or
apply the two component functions and pair. For implication out of a
disjunction, restrict to each injection, or split the supplied disjunction
and use the appropriate component function. Taking $C=\bot$ gives the
disjunction De Morgan law, and currying at $C=\bot$ gives the negated
conjunction identity. From $\neg A$ and $\neg\neg A$ one gets falsehood,
so $\neg A\to\neg\neg\neg A$. Conversely, given $\neg\neg\neg A$, assume
$A$; this assumption gives $\neg\neg A$, contradicting the given triple
negation. Discharge $A$.

Hypothetical syllogism applies $A\to B$ and then $B\to C$ to an assumed
$A$. Contraposition applies a given $A\to B$ under $A$ and contradicts
$\neg B$. Disjunctive syllogism splits $A\vee B$; the $A$ branch contradicts
$\neg A$, and the $B$ branch is immediate. Constructive dilemma splits
$A\vee B$, applies respectively $A\to C$ or $B\to D$, and introduces
$C\vee D$. These instructions specify each proof node and all discharges;
they are fixed finite templates. The structural substitution theorem
makes the templates uniform over formula codes.
\end{proof}

The Boolean equivalences $A\to B\leftrightarrow\neg A\vee B$ and
$\neg(A\wedge B)\leftrightarrow\neg A\vee\neg B$ are intentionally absent
from this intuitionistic list. Their classical validity will be proved
by truth tables, and finite Kripke countermodels below explain why the
unrestricted intuitionistic versions fail.

\subsection{Curry--Howard and typed SK combinators at the finite level}

For the implication-only fragment, use terms $x$, $tu$, and $\lambda x.t$.
A typing context assigns formulas to finitely many variables. The rules
are variable lookup, application from $A\to B$ and $A$, and abstraction
from a derivation of $B$ with one extra assumption $x:A$.

\begin{theorem}[Finite Curry--Howard correspondence and substitution]\label{thm:curry-howard}
Annotated typing derivations and implication-only natural deductions
translate into each other primitive recursively, preserving the conclusion
type. If $\Gamma,x:A\vdash t:B$ and $\Gamma\vdash u:A$, then
$\Gamma\vdash t[u/x]:B$. Every single beta-reduction preserves typing.
\end{theorem}
\begin{proof}
Assign a term variable to each context slot or discharged assumption class,
using that same variable at every corresponding assumption leaf; record application
at implication elimination, and record abstraction at implication
introduction. Conversely erase the term decorations, retaining the rule
tree. These are direct recursive operations. Substitution is the proof
cut transformation from Theorem~\ref{thm:structural}; use fresh numerical
variable labels or de Bruijn indices under binders. The contraction
$(\lambda x.t)u\longrightarrow t[u/x]$ is well typed by this substitution
lemma. For a contraction inside a larger term, rebuild the typing rules
on the path to its root. This proves single-step subject reduction.
We make no claim here that every arbitrary normalization strategy has a
particular primitive recursive running time.
\end{proof}

Let K and S have the types displayed in the Hilbert schemes, read as
curried function types. Put $\mathrm I=\mathrm{SKK}$, with the appropriate
instances of their type schemes. For a lambda-free applicative term define bracket
abstraction by
\[
 [x]x=\mathrm I,\qquad [x]t=\mathrm Kt\ (x\text{ absent from }t),\qquad
 [x](tu)=\mathrm S([x]t)([x]u)
\]
in the remaining application case. To compile a term containing lambdas,
eliminate inner lambdas first. The reduction identity in the next theorem
is stated for the resulting lambda-free applicative body; for an original
body containing lambdas its target is the correspondingly compiled substitution.

\begin{theorem}[Typed bracket abstraction]\label{thm:sk}
The bracket algorithm is primitive recursive and preserves typing:
$\Gamma,x:A\vdash t:B$ gives $\Gamma\vdash[x]t:A\to B$.
Under the reduction rules, closed under all application contexts,
$\mathrm Ktu\to t$ and $\mathrm Sfgx\to fx(gx)$,
the term $([x]t)u$ reduces in a constructively produced finite sequence
to $t[u/x]$. Thus every implication-only proof has a typed SK certificate.
For lambda-free $t$ one has $|[x]t|\le5|t|$, and the constructed reduction
uses at most $2|t|$ K/S contractions. Compiling an arbitrary typed lambda
term of $s$ nodes therefore yields a combinator term of at most $5^s$ nodes.
\end{theorem}
\begin{proof}
The variable case uses the identity type. In the absent-variable case,
K applied to $t:B$ has type $A\to B$. In the application case, write
$t:C\to B$ and $u:C$; the recursively constructed abstractions have
types $A\to(C\to B)$ and $A\to C$, so S produces $A\to B$.
For the reduction claim, $\mathrm Iu=\mathrm{SKK}u$ reduces first to
$\mathrm Ku(\mathrm Ku)$ and then to $u$; $\mathrm Ktu$ reduces to $t$;
and the S case reduces to two smaller abstraction applications, to which
the induction hypotheses apply. Concatenate these finite reduction lists.
The size bound follows because the identity case has five nodes, the K
case adds two nodes, and the S case adds three nodes to the two recursive
outputs. The reduction count is two in the identity case, one in the K
case, and one plus the two recursive counts in the S case. Induction on
the source lambda term then gives the exponential compilation bound.
All intermediate terms and proofs are explicitly constructed by recursion.
No assertion about arbitrary, untyped combinator normalization is used.
\end{proof}
